% TOP_PROBLEM: {"id": 3900015, "title": "Packing reciprocal rectangles in a square", "category_id": 6, "category": {"id": 6, "name": "geometry", "display_name": "Geometry", "description": "Euclidean and non-Euclidean geometry, geometric structures.", "slug": "geometry", "order_index": 6, "created_at": "2026-07-31T15:26:25.671Z"}, "status": "partially_solved"}
% FIRST_POSTED: null
% ENTRY_KIND: statement_only
% Imported from: batch13.tex; UnsolvedMath ID 3900015
\documentclass[11pt]{book}
\usepackage{fontspec}
\setmainfont{Latin Modern Roman}
\setsansfont{Latin Modern Sans}
\usepackage{amsmath,amssymb,mathtools}
\usepackage{unicode-math}
\setmathfont{Latin Modern Math}
\usepackage{xurl}
\usepackage[hidelinks]{hyperref}
\newcommand{\sref}[2]{\hyperlink{source:#1:#2}{[S#2]}}
\newcommand{\cataloguescope}[1]{\paragraph{Mathematical question.}#1}
\begin{document}
% UnsolvedMath ID 3900015; source catalogue code AMR-038-0015
\setcounter{section}{3900014}
\section{Packing reciprocal rectangles in a square}\label{3900015}
{\small\sffamily UnsolvedMath ID 3900015 · Geometry}
\subsection{Definitions and mathematical statement}

\paragraph{Shared notation.}$\mathbb N=\{1,2,\ldots\}$, and $\mathbb Z,\mathbb Q,\mathbb R,\mathbb C$ denote the integers, rationals, reals and complex numbers. Any other conventions are specified locally.


For $k\in\mathbb N$, let
\[
 R_k=[0,1/k]\times[0,1/(k+1)]\subset\mathbb R^2,
 \qquad Q=[0,1]^2.
\]
A packing in $Q$ consists of congruent copies $P_k$ of all the $R_k$ such that $P_k\subseteq Q$ and $\operatorname{int}(P_i)\cap\operatorname{int}(P_j)=\varnothing$ whenever $i\ne j$. Thus contacts between boundaries are allowed. Equivalently, each copy has the form $P_k=v_k+U_kR_k$, with $v_k\in\mathbb R^2$ and $U_k\in SO(2)$ a planar rotation. The question concerns one simultaneous packing of the entire infinite sequence.
Since
\[
 \sum_{k=1}^{\infty}\operatorname{area}(R_k)
 =\sum_{k=1}^{\infty}\left(\frac1k-\frac1{k+1}\right)=1,
\]
any such packing leaves an uncovered set of planar area zero; a pointwise tiling is not required. Packings of large finite prefixes do not by themselves specify an infinite packing. \sref{3900015}{2}
\paragraph{Statement recorded in the supplied source.}
Do there exist vectors $v_k$ and rotations $U_k\in SO(2)$ for every $k\in\mathbb N$ such that $v_k+U_kR_k\subseteq Q$ and the interiors of these copies are pairwise disjoint? \sref{3900015}{2}
\subsection{Short English statement}
Can one copy of every $1/k$ by $1/(k+1)$ rectangle, starting with $k=1$, fit in a single unit square without overlapping interiors? The total area is exactly one, so a packing would cover the square except for a set of area zero.
\subsection{Sources}
\begin{description}
\item[\hypertarget{source:3900015:1}{S1}] ULAM.AI and the UnsolvedMath Contributors, \emph{UnsolvedMath: A Curated Collection of Open Mathematics Problems}, record 3900015 (AMR-038-0015). CC BY 4.0. \url{https://huggingface.co/datasets/ulamai/UnsolvedMath}.
\item[\hypertarget{source:3900015:2}{S2}] Mingliang Zhu and Antal Jo\'os, \emph{Packing $1.35\cdot10^{11}$ rectangles into a unit square} (2022), arXiv:2211.10356. Section~1, pp.~1--2: definition of packing and complete Meir--Moser rectangle sequence; section~2, finite-prefix results; section~3, distinction from the infinite question. \url{https://arxiv.org/abs/2211.10356}.
\end{description}
\label{end:3900015}
\end{document}
